\section*{Capítulo \ref{ch:b3:homogeneous-catalysis} --- Catálisis homogénea en la industria}

\begin{solution}{exo:b3:homogeneous-catalysis:1}
$\ce{RhCl(PPh3)3}$: 16, Rh(I) $d^8$. $\ce{RhCl(PPh3)2}$: 14, Rh(I). $\ce{RhClH2(PPh3)2}$: 16, Rh(III) $d^6$.
Complejo de alqueno: 18, Rh(III). Hidruro de alquilo: 16, Rh(III). La eliminación reductora
devuelve a 14, Rh(I).
\end{solution}

\begin{solution}{exo:b3:homogeneous-catalysis:2}
Butanal, $\ce{CH3CH2CH2CHO}$ (lineal), y 2-metilpropanal, $\ce{(CH3)2CHCHO}$ (ramificado).
\end{solution}

\begin{solution}{exo:b3:homogeneous-catalysis:3}
Ciclopent-3-eno-1,1-dicarboxilato de dietilo, un anillo de cinco miembros, con pérdida de
eteno: los dos grupos $\ce{CH2}$ terminales se van como $\ce{C2H4}$ y los dos carbonos CH internos
se unen.
\end{solution}

\begin{solution}{exo:b3:homogeneous-catalysis:4}
(\textit{E})-3-fenilprop-2-enoato de metilo (cinamato de metilo): el arilo se une al
carbono terminal, y la eliminación de hidruro en $\beta$, sin, desde el confórmero más estable
da el alqueno \textit{E}.
\end{solution}

\begin{solution}{exo:b3:homogeneous-catalysis:5}
La adición oxidante es lenta y el $\ce{[Rh(CO)2I2]-}$ es el estado de reposo, de modo que
$v = k[\mathrm{Rh}]_{\mathrm{tot}}[\ce{CH3I}]$. El metanol solo regenera el yoduro de metilo por una reacción rápida con HI;
el yoduro total cargado fija $[\ce{CH3I}]$, de modo que la velocidad no depende de la concentración de
metanol hasta que queda demasiado poco metanol para volver a convertir el HI en
yoduro de metilo.
\end{solution}

\begin{solution}{exo:b3:homogeneous-catalysis:6}
Ácido fenilborónico con 4-bromotolueno, o ácido 4-metilfenilborónico con
bromobenceno. Ambas funcionan; la primera usa el ácido borónico más barato y un bromuro de arilo
estable y corriente.
\end{solution}

\begin{solution}{exo:b3:homogeneous-catalysis:7}
Simetría $C_2$: isotáctico (ambas posiciones seleccionan la misma cara). Simetría $C_s$:
sindiotáctico (las posiciones especulares alternan). $\ce{Cp2ZrCl2}$ sin puente: atáctico (posiciones
aquirales).
\end{solution}

\begin{solution}{exo:b3:homogeneous-catalysis:8}
$\mathrm{TON} = 0.98/0.0010 = 980$; de media, $\mathrm{TOF} = 980/\qty{2.0}{h} = \qty{490}{h^{-1}}$.
\end{solution}

\begin{solution}{exo:b3:homogeneous-catalysis:9}
La unidad repetitiva es $\ce{-[CH=CH-C5H8]-}$, con los dos grupos CH en los carbonos 1 y 3 de un
anillo de ciclopentano (poli(1,3-ciclopentilenvinileno)). La liberación de la
tensión del anillo bicíclico hace $\Delta_rH^\circ$ muy negativa, lo que compensa con creces la
pérdida de entropía de traslación; no se forma ningún gas, de modo que la fuerza impulsora es
entálpica.
\end{solution}

\begin{solution}{exo:b3:homogeneous-catalysis:10}
$\bar n_n = 1/(1 - p) = 2000$; $M_w/M_n = 1 + p = 1.9995 \approx 2.0$. Un \omterm{def:b3:homogeneous-catalysis:ziegler-natta}{catalizador de Ziegler--Natta} tiene
varios tipos de sitios, cada uno con su propio $p$; la suma de varias distribuciones de
Schulz--Flory de medias distintas es más ancha (dispersidad a menudo de 4 a 8).
\end{solution}

\begin{solution}{exo:b3:homogeneous-catalysis:11}
A baja presión, $K_1K_2[\ce{H2}] \ll [\mathrm L] + K_1$: $v \approx kK_1K_2[\mathrm{Rh}]_{\mathrm{tot}}[\text{alqueno}][\ce{H2}]/([\mathrm L] + K_1)$,
de primer orden respecto del $\ce{H2}$. A alta presión domina el término $K_1K_2[\ce{H2}]$: $v \to k[\mathrm{Rh}]_{\mathrm{tot}}[\text{alqueno}]$,
independiente del $\ce{H2}$. $[\mathrm L]$ solo aparece en el denominador: la fosfina añadida reduce $v$
al devolver el rodio al precatalizador saturado.
\end{solution}

\begin{solution}{exo:b3:homogeneous-catalysis:12}
La inserción que sitúa el rodio sobre el carbono interno forma un alquilo secundario
junto a un metal congestionado; las fosfinas voluminosas (y un exceso de ellas, que mantiene dos
sobre el metal) hacen ese estado de transición más desfavorable que el que da el alquilo
primario, que conduce al aldehído lineal. El butanal lineal es el que se convierte,
por condensación aldólica e hidrogenación, en el alcohol ramificado de $C_8$ usado en los
plastificantes; el 2-metilpropanal es un subproducto de menor valor.
\end{solution}

\begin{solution}{pb:b3:homogeneous-catalysis:1}
\textbf{1.} Ir(I), $d^8$: $9 + 4 + 2 + 1 = 16$ electrones (recuento neutro, con la carga).

\textbf{2.} $\ce{[CH3Ir(CO)2I3]-}$, 18 electrones, Ir(III).

\textbf{3.} $\ce{[CH3Ir(CO)3I2]}$, 18 electrones, Ir(III).

\textbf{4.} $\ce{[CH3COIr(CO)2I2]}$, 16 electrones, Ir(III).

\textbf{5.} El $\ce{[CH3COIr(CO)2I3]-}$ (18) elimina $\ce{CH3COI}$ y regenera el $\ce{[Ir(CO)2I2]-}$.

\textbf{6.} $\ce{CH3COI + H2O -> CH3COOH + HI}$; $\ce{CH3OH + HI -> CH3I + H2O}$.

\textbf{7.} $\ce{CH3OH + CO -> CH3COOH}$.

\textbf{8.} Con el iridio, la inserción migratoria es lenta; la adición oxidante del
yoduro de metilo, determinante de la velocidad con el rodio, es mucho más rápida sobre el iridio, más
rico en electrones.

\textbf{9.} Retiran un yoduro del $\ce{[CH3Ir(CO)2I3]-}$, lo que da el tricarbonilo neutro, en el que
el metal retrodona menos y el metilo migra más deprisa hacia un CO más
electrófilo.

\textbf{10.} Una dependencia inversa de la concentración de yoduro libre: el yoduro desplaza
el preequilibrio de vuelta hacia el complejo aniónico lento.

\textbf{11.} El metanol solo entra después de la etapa determinante de la velocidad, convertido en
yoduro de metilo por una reacción rápida con HI.

\textbf{12.} $M(\ce{CH3COOH}) = \qty{60.1}{g/mol}$: $n = \qty{5.0e11}{g}/\qty{60.1}{g/mol} = \qty{8.3e9}{mol}$ al año.

\textbf{13.} $\qty{8.32e9}{mol}/\qty{8000}{h} = \qty{1.04e6}{mol/h}$.

\textbf{14.} $\qty{2.0e5}{L} \times \qty{5.0e-3}{mol/L} = \qty{1.0e3}{mol}$ de iridio, $\qty{1.0e3}{mol} \times \qty{192.2}{g/mol} \approx \qty{190}{kg}$.

\textbf{15.} $\mathrm{TOF} = \qty{1.04e6}{mol/h}/\qty{1.0e3}{mol} = \qty{1.0e3}{h^{-1}}$, unos \qty{0.29}{s^{-1}}.

\textbf{16.} $\mathrm{TON} = \qty{8.3e9}{mol}/\qty{1.0e3}{mol} = \num{8.3e6}$ al año.

\textbf{17.} $\qty{5.0e8}{kg}/(\qty{8000}{h} \times \qty{200}{m^3}) \approx \qty{310}{kg.m^{-3}.h^{-1}}$.

\textbf{18.} $\ce{CO + H2O -> CO2 + H2}$.

\textbf{19.} $1/0.94 = \qty{1.064}{mol}$ de CO por mol de ácido acético.

\textbf{20.} $1.064 - 1 = \qty{0.064}{mol}$ de \ce{CO2} por mol de ácido acético.

\textbf{21.} Una tonelada son $\qty{1.0e6}{g}/\qty{60.1}{g/mol} = \qty{1.66e4}{mol}$; $\qty{1.66e4}{mol} \times 0.0638 \times \qty{44.0}{g/mol} \approx \qty{47}{kg}$ de \ce{CO2}.

\textbf{22.} El agua es un reactivo de la reacción de desplazamiento: menos agua la ralentiza y ahorra la energía
del secado del ácido. El agua sigue siendo necesaria para hidrolizar el yoduro de acetilo y para mantener el
catalizador activo y disuelto; con el rodio, poca agua hace precipitar el yoduro de
rodio.

\textbf{23.} Hidrógeno: se acumula en el gas, rebaja la presión parcial de CO
(de modo que hay que purgar gas, con pérdida de CO) y puede hidrogenar intermedios y dar
subproductos como el ácido propiónico.

\textbf{24.} El catalizador de iridio efectúa unos $\qty{1.0e3}{h^{-1}}$ ciclos: cada átomo de iridio fabrica
mil moléculas de ácido acético por hora, unos ocho millones al año.
\end{solution}
